\documentclass[11pt]{article}
\usepackage[T1]{fontenc}
\usepackage{amsmath,amssymb,amsthm,mathtools}
\usepackage[margin=1in]{geometry}
\usepackage[colorlinks=true,allcolors=blue]{hyperref}
\setlength{\emergencystretch}{2em}
\newtheorem{theorem}{Theorem}[section]
\newtheorem{proposition}[theorem]{Proposition}
\newtheorem{lemma}[theorem]{Lemma}
\newtheorem{corollary}[theorem]{Corollary}
\newtheorem{problem}[theorem]{Problem}
\theoremstyle{definition}
\newtheorem{definition}[theorem]{Definition}
\newtheorem{remark}[theorem]{Remark}
\newcommand{\Rep}{\operatorname{Rep}}
\newcommand{\Hom}{\operatorname{Hom}}
\newcommand{\End}{\operatorname{End}}
\newcommand{\Ext}{\operatorname{Ext}}
\newcommand{\Db}{D^b}
\title{Cyclic symmetry and failure of block-permutation invariance in Dexter intervals}
\author{YINGFENG JIANG}
\date{7 October 2026}
\begin{document}
\maketitle
\begin{abstract}
Over every field, cyclic rotation of the blocks of a Dyck word
preserves the bounded derived category of its Dexter interval $J(w)$.
Arbitrary block permutation does not: an exact computation gives
two 288-element intervals with distinct integral Euler polynomials.
For the three blocks $110100$, $11011000$, and $1100$, their six
permutations consequently give exactly two derived equivalence classes,
namely the two cyclic orbits. Boundary attachment explains the distinction.
The lower truncations $L_{uv}$ and $L_{vu}$ for
$u=110100$, $v=11011000$ are derived equivalent, but no enhanced
equivalence between them intertwines their
boundary pullbacks, even after an arbitrary enhanced autoequivalence
of the base. We prove this using a relative extension criterion and
the counterexample. Exchanges involving the elementary block and
transfers between pyramid blocks provide positive supplementary families.
\end{abstract}

\section{Introduction}
\label{sec:introduction}
A Dyck word is a word in $1,0$ with nonnegative prefix heights and total
height zero. Its blocks are its nonempty factors between consecutive
returns to height zero. Write $b(x)$ for the number of blocks of $x$.
We use the Dexter order of \cite{Chapoton}. If a primitive Dyck subpath
$s$ is preceded by a zero and is followed by a one or by the end of
the word, it is movable. A cover slides $s$ to the left across any
positive number of consecutive immediately preceding zeros. Such a
subpath may be based at positive height; it need not be a whole block.

Let $I(w)=[w_{\min},w]$ and
\[
 J(w)=[w_{\min},(1,w,1,0,0)].
\]
For a finite poset $P$, $\Rep_k P$ denotes finite-dimensional covariant
representations, and all equivalences are $k$-linear. The following
assertion is the block-permutation subproblem explicitly proposed in
\cite[Section 9, Conjecture 9.1]{Chapoton}.
\begin{problem}[Block-permutation assertion, disproved below]\label{prob:swap}
For every field $k$, every pair of nonempty indecomposable Dyck blocks
$u,v$, and possibly empty block strings $A,B$, does one have
\[
 \Db(\Rep_k J(AuvB))\simeq\Db(\Rep_k J(AvuB))?
\]
\end{problem}
The positive symmetry is cyclic rotation: for all nonempty Dyck words
$u,v$ and every field $k$, Theorem~\ref{thm:rotation} gives
\[
 \Db(\Rep_k J(uv))\simeq\Db(\Rep_k J(vu)).
\]
In contrast, with $u=110100$, $v=11011000$, and $b=1100$,
Theorem~\ref{thm:counterexample} separates $J(uvb)$ from $J(vub)$
by their integral Euler polynomials. Both intervals have 288 points.
Corollary~\ref{cor:two-cyclic-orbits} combines the
two results into exactly two equivalence classes for one specified
three-block multiset. No general classification by cyclic order is claimed.

The distinction is visible in the boundary data. Theorem~\ref{thm:lower-rotation}
identifies the abstract derived categories of the lower truncations.
The relative extension criterion, Proposition~\ref{prop:relative-extension},
then shows that for the counterexample pair their boundary embeddings
cannot be intertwined by any enhanced equivalence and any enhanced
base autoequivalence (Corollary~\ref{cor:absolute-boundary-obstruction}).
Thus the boundary-decorated category contains information absent from
its abstract derived type. Elementary exchanges and pyramid transfers
are supplementary positive results, not steps toward the disproved
universal assertion.

The product decomposition of $I(w)$, the chain decomposition, and the
canonical shape cut are inputs from \cite[Proposition 4.4 and Section 5.2]{Chapoton}.
For general finite posets, Ladkani proved that the ordinal sums
$X\oplus Y$ and $Y\oplus X$ are derived equivalent
\cite[Corollary 4.15]{LadkaniSheaves}; associativity then gives cyclic
rotation of any ordinal sum. His Example 4.20 shows that a permutation
of three summands can fail to preserve derived equivalence.
Thus cyclic symmetry together with permutation failure has a precedent
for ordinal sums. The singleton flip-flop mechanism used in
Appendix~\ref{app:shape} is also known
\cite[Corollary 1.3]{LadkaniPosets}, and the triangular-algebra tilting
method has the background of \cite[Theorems 4.5 and 4.9]{LadkaniTriangular}.

Here the objects are the specific Dexter intervals $J(w)$, whose
mixed relations depend on their projecting boundary coordinates.
The calculations below identify the cross-chain rules, the target
mixed bimodules and their multiplication, and the natural boundary
comparisons needed for rotation and pyramid transfer. The explicit
block-permutation counterexample concerns these intervals.
The ordinal-sum results alone do not identify the required Dexter
objects and boundary data; their relation is not assumed in our proofs.
Appendix calculations use different examples to explain limitations
of particular cuts and restricted base corrections.

The paper follows this distinction. Section~\ref{sec:structure} gives
the chain coordinates and the marked-boundary algebra.
Section~\ref{sec:rotation} proves cyclic symmetry; Section~\ref{sec:failure}
then supplies the counterexample and the two cyclic orbits.
Section~\ref{sec:boundary} proves the abstract lower-truncation
equivalence, the relative extension criterion, and their incompatible
boundary consequence, in that order. Section~\ref{sec:positive}
records the positive supplementary families. Computational details
and the longer mutation argument are placed in the appendices.


\section{Chain coordinates and the marked-boundary algebra}
\label{sec:structure}
\subsection{The fibres and their cross-chain arrows}
For $x=x_1\cdots x_b$ in blocks, put
\begin{align*}
 e_i(x)&=(1,x_1,\ldots,x_i,0,x_{i+1},\ldots,x_b,1,0)
                       &&(0\leq i\leq b),\\
 e_{b+1}(x)&=(1,x,1,0,0).
\end{align*}
We call the first $b+1$ points lower-shape points and the last point
a top-shape point. This terminology avoids confusing shape labels
with the prefix and suffix $A,B$.

\begin{lemma}\label{lem:fibres}
The elements of $J(w)$ are precisely $e_i(x)$ with $x\leq w$ and
$0\leq i\leq b(x)+1$. Each such chain is complete, and the map
$q(e_i(x))=x$ is order preserving. The top-shape subposet is $I(w)$.
\end{lemma}
\begin{proof}
The shape subposet of \cite[Lemmas 5.3--5.4]{Chapoton} is a lower
ideal partitioned by these chains. A cross-chain cover projects to a
lower cover by Proposition 5.5 of that paper. Thus $q$ is order
preserving. A point below the top $e_{b(w)+1}(w)$ must project below
$w$. Conversely the top-shape embedding preserves order, so the
entire chain above $x\leq w$ lies below that top.
\end{proof}

A lower cover $a:x\prec y$ has one of two types. An internal cover
slides a subpath based at positive height and leaves the block count
unchanged. A merging cover slides a whole block $x_\ell$, with
$2\leq\ell\leq b(x)$, into its preceding block. Then $b(y)=b(x)-1$;
we call $\ell$ the merging position. These exhaust the covers:
a slide cannot cross a one, so a zero-height slide joins exactly two
adjacent blocks.

\begin{proposition}[Cross-chain cover rule]\label{prop:covers}
All covers of $J(w)$ are the following, with $x,y\leq w$.
\begin{enumerate}
\item The vertical covers $e_i(x)\prec e_{i+1}(x)$.
\item For an internal cover $x\prec y$, the covers
$e_i(x)\prec e_i(y)$ for $0\leq i\leq b(x)+1$.
\item For a merging cover $x\prec y$ at position $\ell$, the covers
\[
 e_i(x)\prec e_i(y)\quad(0\leq i<\ell),\qquad
 e_i(x)\prec e_{i-1}(y)\quad(\ell<i\leq b(x)+1).
\]
There is no cross-chain cover over this lower cover with source
$e_\ell(x)$.
\end{enumerate}
\end{proposition}
\begin{proof}
Write a lower-shape point as $(1,p,0,q,1,0)$, where $x=pq$ and
$p$ consists of the first $i$ blocks. Call the zero between $p$ and
$q$ the inserted zero.
An internal movable subpath has positive baseline in $x$.
It cannot end at a zero-height cut, and its preceding zero run cannot
cross such a cut. Inserting the wrapper therefore preserves its
movability and its slide, on either side of the cut. The index stays $i$.

Now slide the whole block $x_\ell$ across $t$ zeros in $x$.
If $i<\ell-1$, the move occurs in $q$ and the cut index is unchanged.
If $i=\ell-1$, slide it across $t+1$ zeros in the wrapped word:
the extra zero is the inserted zero. The resulting wrapped prefix
ends after the merged block, at index $\ell-1$ of $y$.
If $i>\ell$, the move lies in $p$, and $x_\ell$ is followed there by
a one. It remains movable; merging reduces the prefix block count
by one. If $i=\ell$, however, $x_\ell$ is followed by the inserted
zero. It is not movable. This is precisely the excluded case.
The top-shape case is the unchanged lower slide inside the wrapper,
with top index reduced by one when a merge occurs.

For completeness, consider any slide that remains in the shape
subposet. A slide of the first block of $q$ across just the inserted
zero is a vertical cover. Sliding it farther gives exactly the
merging case just described. All remaining slides away from the
terminal $10$ project to an internal or merging lower cover.
A slide of that terminal $10$ stays in the shape subposet only when
$q$ is empty and it crosses just the inserted zero; this is the
last vertical cover. If $q$ is nonempty, such a slide either leaves
an earlier return to height zero or produces too many terminal zeros,
and is neither of the two shapes. A longer terminal slide when
$q$ is empty also has too many terminal zeros. These observations
exhaust the movable subpaths. Since the shape subposet is a lower
ideal, its induced covers are precisely the ambient Dexter covers
with both endpoints in it. Restricting to $J(w)$ introduces no new
covers. The stated list is therefore complete.
\end{proof}

The excluded position cannot be repaired by pretending that every
cut simply follows its image in the merged word. The correct formula
for comparability allows vertical motion first.
For each lower cover $a:x\prec y$ define a monotone map between
index chains by
\[
 f_a(i)=
 \begin{cases}
 i,&a\text{ internal},\\
 i-\mathbf1_{i>\ell},&a\text{ a merge at }\ell.
 \end{cases}
\]
The domains include the top index $b(x)+1$.

\begin{corollary}[Full comparability]\label{cor:threshold}
For $x\leq y$ let
\[
 m_{x,y}(i)=\min_{\pi:x\leadsto y} f_\pi(i),
 \qquad f_\pi=f_{a_r}\circ\cdots\circ f_{a_1},
\]
where $\pi$ runs through saturated lower paths; for $x=y$ use the
empty path. Then
\[
 e_i(x)\leq e_j(y)\quad\Longleftrightarrow\quad
 x\leq y\ \hbox{and}\ j\geq m_{x,y}(i).
\]
\end{corollary}
\begin{proof}
Over one lower cover, $f_a(i)$ is the smallest reachable target
index. The cover rule proves this directly except at $i=\ell$.
There one first moves vertically to $\ell+1$ and then crosses to
index $\ell$; if $\ell=b(x)$ the intermediate vertex is the top.
No smaller target is reachable by the listed covers. Monotonicity
shows that waiting longer on a fibre cannot improve this minimum.
Any upper path projects, after deleting vertical steps, to a lower
saturated path. Its final index is at least the corresponding
composite $f_\pi(i)$. Conversely each such composite can be attained
by lifting its covers in order; any larger final index is obtained
by vertical motion. There are finitely many paths, so the minimum
is attained.
\end{proof}
\begin{remark}
This formula takes a minimum over paths. It does not assume that
different lower paths give the same map of chains, and does not
assert a Grothendieck construction with path-independent transition
functors.
\end{remark}


\subsection{Incidence algebras and exterior factors}
Fix the category-algebra convention in which a basis element
$a_{ts}$ represents $s\leq t$, and $a_{ut}a_{ts}=a_{us}$.
Covariant representations are left modules. By
Corollary~\ref{cor:threshold}, the algebra has basis $a_{ts}$ for
exactly the comparable chain points, and this composition law.
Equivalently it is the path algebra of the quiver in
Proposition~\ref{prop:covers}, modulo the differences of all parallel
paths. Indeed each comparable pair has a path, these relations
identify all paths with those endpoints, and relations with distinct
endpoints do not interact. The quotient therefore has precisely the
claimed basis, including the identity paths.

The lower product decomposition gives
\[
 I(AuvB)=I(A)\times I(u)\times I(v)\times I(B),
\]
where the empty word contributes a one-point factor. This is an
actual product of posets, not a claim about $J(AuvB)$.
\begin{corollary}[Context offsets]\label{cor:context}
Suppose a cover occurs in one factor of this product, and let $c$
be the total number of blocks in the current preceding factors.
An internal move uses the identity index rule. A merging move at
local position $h$ uses the rule of Proposition~\ref{prop:covers}
with global merging position $\ell=c+h$. The incidence-algebra
presentation above includes all arrows and mixed compositions for
arbitrary $A,B$.
\end{corollary}
\begin{proof}
The block boundaries of the upper word must remain returns to zero
throughout its lower interval. Thus a cover changes exactly one
factor. Counting the blocks before that factor gives $\ell=c+h$;
the complete cover rule and the parallel-path presentation apply.
\end{proof}
The offset $c$ depends on the lower element, not only on the fixed
upper word. Swapping product factors does not by itself transport
these cover rules. Corollary~\ref{cor:context} supplies the original
attachments. Transport through a derived equivalence additionally
requires the boundary comparison in Section~\ref{sec:boundary}.


\subsection{Derived conventions and tilting}
For a vertex $x$ of a finite poset, write
$P_x=k\operatorname{Hom}(x,-)$. Thus
$\operatorname{Hom}(P_x,P_y)=k$ exactly when $y\leq x$.
For a tilting object indexed by a target poset, the algebra to identify
is its opposite endomorphism algebra. We use the following specified
finite-algebra functors. For a finite poset $X$, write
$\mathcal D_X=\Db(\Rep_k X)$ and denote its Serre functor by
$\mathbb S_X$. It is derived tensor with the $k$-dual of $kX$;
finite global dimension makes it an autoequivalence. Derived
Kan extensions along $q:X\to S$ are denoted
\[
 \ell_X\dashv Q_X\dashv r_X,\qquad Q_X=q^*.
\]
They can be computed by finite poset bar resolutions. Their Serre
adjunction identity is
\begin{equation}\label{eq:serre-adjoint}
 r_X\mathbb S_X\simeq\mathbb S_S\ell_X.
\end{equation}
Indeed Serre duality applied to
$\Hom_X(Q_X V,\mathbb S_X E)$ identifies it with
$\Hom_S(V,\mathbb S_S\ell_X E)$, naturally in both variables.
In particular $\ell_X(P_a)=P_{q(a)}$. External tensor products
and derived Hom below are over the field $k$.
We use the standard differential graded enhancements by perfect
module complexes. Enhanced equivalences and functor comparisons
mean equivalences and isomorphisms of the corresponding derived
bimodules; in particular they admit external tensor products and
transport both actions of a gluing bimodule. All functors constructed
below are represented in this framework. The negative boundary
statement quantifies over this entire class of enhanced equivalences.

\begin{lemma}[Tilting convention]\label{lem:tilting-principle}
Let $T$ generate $\mathcal D_X$, satisfy
$\operatorname{Hom}(T,T[m])=0$ for $m\ne0$, and have
$\operatorname{End}(T)^{\mathrm{op}}\cong kY$ for a finite poset $Y$,
respecting the specified vertex idempotents.
Then there is an enhanced equivalence $\mathcal D_Y\to\mathcal D_X$
sending the vertex projectives to the specified summands of $T$.
\end{lemma}
\begin{proof}
A projective-complex representative of $T$ gives a
differential graded endomorphism algebra with cohomology only in
degree zero. Its truncation provides a quasi-isomorphism zigzag to
$\End(T)$. The Hom functor is an equivalence on the thick category
generated by $T$: it is so on the generator and hence on finite
cones and summands, whose image is the perfect module category.
Generation and finite global dimension of the identified target
poset algebra give the claimed bounded equivalence.
\end{proof}


\subsection{Cutting at a word boundary}
For a nonempty Dyck word $z$, put $S_z=I(z)$ and define two full
subposets of $J(z)$:
\[
 L_z=\{e_i(x):x\leq z,\ 0\leq i\leq b(x)\},\qquad
 R_z=\{e_i(x):x\leq z,\ 1\leq i\leq b(x)+1\}.
\]
Thus $L_z$ removes the top-shape layer, whereas $R_z$ removes the
bottom layer of every chain. Both carry the restriction $q_z$ of
$q$. A poset subscript on a derived functor specifies its actual finite
representation category.

\begin{proposition}[The marked-boundary algebra]\label{prop:boundary-algebra}
For nonempty Dyck words $u,v$, the poset $J(uv)$ is a two-piece
poset with lower ideal $D=L_u\times S_v$ and upper filter
$U=S_u\times R_v$. Its mixed relations are exactly
\begin{equation}\label{eq:boundary-mixed}
 (d,y)\leq(x,r)\quad\Longleftrightarrow\quad
 q_u(d)\leq x\ \text{and}\ y\leq q_v(r).
\end{equation}
These are relations of the actual poset, including their incidence
compositions.
\end{proposition}
\begin{proof}
Write a lower element as $xy$, with $x\leq u$, $y\leq v$, and put
$a=b(x)$. Partition its chain at $i=a$: the point $e_i(xy)$ belongs
to $D$ if $i\leq a$, with coordinates $(e_i(x),y)$; otherwise it
belongs to $U$, with coordinates $(x,e_{i-a}(y))$.
For a lower cover inside $u$, Proposition~\ref{prop:covers} restricts
on $D$ to the $L_u$ rule and acts on $U$ without changing the
$R_v$ index: both $i$ and $a$ decrease by one on a merging move.
For a lower cover inside $v$, it restricts on $U$ to the $R_v$ rule
and leaves the $L_u$ index unchanged on $D$. A crossing can leave
$D$ for $U$, but none can go from $U$ to $D$. This proves the
claimed induced subposets and ideal/filter properties.

The section $h(x,y)=e_{b(x)+1}(xy)$ is monotone: merges inside $u$
reduce its index by one, and merges inside $v$ preserve it. It is
the first point of $U$ in its fibre. Every $d$ lies below the
section in its own fibre, and $h(x,y)$ lies below every upper-fibre
point over $(x',y')\geq(x,y)$. Monotonicity of $q$ proves the
converse implication. This gives \eqref{eq:boundary-mixed}.
The composition statement follows from the incidence basis, not
from an independent choice of mixed scalars.
\end{proof}



\section{Cyclic rotation}
\label{sec:rotation}
\subsection{The two truncations and their relative equivalence}
\begin{lemma}[An explicit truncation equivalence]\label{lem:truncation}
There is an equivalence $F_z:\mathcal D_{L_z}\to\mathcal D_{R_z}$
with a natural isomorphism
\begin{equation}\label{eq:relative-F}
 F_z Q_{L_z}\simeq Q_{R_z}\mathbb S_{S_z}.
\end{equation}
Consequently
\begin{align}
 \ell_{L_z}F_z^{-1}&\simeq\mathbb S_{S_z}^{-1}\ell_{R_z},
       \label{eq:F-left}\\
 r_{R_z}F_z\mathbb S_{L_z}&\simeq
             \mathbb S_{S_z}^{2}\ell_{L_z}.
       \label{eq:F-right}
\end{align}
All these identifications are identities of functors, with their
maps, not only identities on object classes.
\end{lemma}
\begin{proof}
Work first in $\mathcal C=\mathcal D_{J(z)}$.
Let $B=Q_{J(z)}:\mathcal D_{S_z}\to\mathcal C$ and let
$Z:\mathcal D_{S_z}\to\mathcal C$ be extension from the top layer.
For projectives they are respectively
\[
 B(P_x)=P_{e_0(x)},\qquad Z(P_x)=P_{e_{b(x)+1}(x)}.
\]
Thus both are fully faithful. The first formula follows because
$e_0(x)\leq e_i(y)$ exactly when $x\leq y$; the second is extension
from an upper filter. They have both adjoints on these bounded
categories. Injectives at top points give the natural identity
\[
 \mathbb S_{J(z)}Z\simeq B\mathbb S_{S_z}.
\]
Indeed the injective at the top over $x$ has value $k$ exactly on
points projecting below $x$, with the usual incidence maps.

Use extension by zero to regard $\mathcal D_{L_z}$ and
$\mathcal D_{R_z}$ as full derived subcategories of $\mathcal C$.
These embeddings are fully faithful: the sets are respectively a
lower ideal and an upper filter, and restriction of their
projective resolutions verifies derived full faithfulness.
They identify with $Z^\perp$ and $B^\perp$, respectively, by
evaluation on the top and bottom projectives. Serre duality and
the displayed identity give
$Z^\perp={}^\perp B$. Here orthogonality includes all shifts.

For $E\in{}^\perp B$, form the actual counit cone
\[
 H_z(E)=\operatorname{Cone}(B B^R E\longrightarrow E).
\]
It lies in $B^\perp$. This defines an equivalence
${}^\perp B\to B^\perp$. To check full faithfulness, apply Hom to
the two counit triangles: $\Hom(E,B(-)[m])=0$ and
$\Hom(B(-),H_z(E')[m])=0$ identify the Hom groups with those before
mutation. An inverse on $G\in B^\perp$ is the fibre of the unit
$G\to B B^L G$; it belongs to ${}^\perp B$, and the same triangles
show that the two operations are inverse. These cones and fibres
are taken in the derived module enhancement.

The exact sequence
$0\to Z(P_x)\to B(P_x)\to Q_{L_z}(P_x)\to0$
shows, naturally on the projective generators of $\mathcal D_{S_z}$,
that
\[
 H_z Q_{L_z}\simeq Z_{R_z}[1],
\]
where $Z_{R_z}$ is the top-layer projective embedding into $R_z$.
In $R_z$ the same injective calculation gives
$\mathbb S_{R_z}Z_{R_z}\simeq Q_{R_z}\mathbb S_{S_z}$.
Notice that $Q_{R_z}(P_x)=P_{e_1(x)}$; the index $1$ is preserved
by every lower cover. Set
\[
 F_z=\mathbb S_{R_z}H_z[-1].
\]
This proves \eqref{eq:relative-F}. Taking left adjoints gives
\eqref{eq:F-left}. Taking right adjoints first gives
$r_{R_z}F_z\simeq\mathbb S_{S_z}r_{L_z}$; composing with
$\mathbb S_{L_z}$ and applying \eqref{eq:serre-adjoint} gives
\eqref{eq:F-right}. Naturality follows from the actual counit,
unit and Serre pairings used throughout.
\end{proof}


\subsection{The rotation tilting object}
The next result treats two factors without exterior strings. The
factors may themselves have several blocks, so it also proves
cyclic rotation of an arbitrary block string.
Its two-piece tilting construction uses the triangular-algebra method
of \cite[Theorems 4.5 and 4.9]{LadkaniTriangular} as background.
We compute the actual mixed bimodule and identify the target
$kJ(vu)$, including multiplication, in the proof below.

\begin{theorem}[Two-factor rotation]\label{thm:rotation}
For all nonempty Dyck words $u,v$ and all fields $k$ there is an
equivalence
\[
 \Db(\Rep_k J(uv))\simeq\Db(\Rep_k J(vu)).
\]
\end{theorem}
\begin{proof}
Use Proposition~\ref{prop:boundary-algebra}, with
$D=L_u\times S_v$ and $U=S_u\times R_v$.
Let $j_0:\mathcal D_D\to\mathcal D_{J(uv)}$ and
$i:\mathcal D_U\to\mathcal D_{J(uv)}$ be extension by zero.
The first is the lower-supported embedding and the second the
upper-supported embedding. For $E\in\mathcal D_D$ and
$V\in\mathcal D_U$, the triangular algebra gives
\begin{align}
 R\Hom(iV,j_0E)&=0,\label{eq:upper-lower-zero}\\
 R\Hom(j_0E[-1],iV)&\simeq R\Hom_U(\Phi E,V),\qquad
 \Phi=\ell_{L_u}\boxtimes Q_{R_v}.
       \label{eq:cross-hom}
\end{align}
Here is a direct justification of the second formula. Derived
extension of the lower projectives gives a triangle
$i\Phi E\to j_!E\to j_0E$; its upper restriction is the actual
mixed bimodule of \eqref{eq:boundary-mixed}. On
$P_d\boxtimes P_y$ that bimodule gives
$P_{q_u(d)}\boxtimes Q_{R_v}P_y$, which proves the formula for
$\Phi$, including its actions, by projective resolutions.
Adjunction gives $R\Hom(j_!E,iV)=0$, and hence
\eqref{eq:cross-hom}. The zero in \eqref{eq:upper-lower-zero}
follows because upper projectives stay projective in the full
poset and have no maps to a lower-supported object.

Index the desired summands by the marked-boundary decomposition
of the target $J(vu)$: its lower vertices are
$(d,x)\in L_v\times S_u$, and its upper vertices are
$(y,r)\in S_v\times R_u$. Define in the source category
\begin{align*}
 T_{d,x}&=i\bigl(\mathbb S_{S_u}^{-1}P_x\boxtimes
                        F_v\mathbb S_{L_v}P_d\bigr),\\
 T_{y,r}&=j_0\bigl(F_u^{-1}P_r\boxtimes
                        \mathbb S_{S_v}^{2}P_y\bigr)[-1].
\end{align*}
This is a uniform construction by finite derived tensor products,
Kan extensions and the counit cones of Lemma~\ref{lem:truncation}.
It is not specified by assuming the required endomorphism algebra.
Each factor is a bounded complex over a finite poset algebra.

Inside each family, derived full faithfulness and the equivalences
$F_z$ and $\mathbb S_X$ identify all Hom complexes with those of the
corresponding target projectives. The Hom complexes from the first
family to the second vanish by \eqref{eq:upper-lower-zero}.
For the remaining direction, \eqref{eq:cross-hom}, the tensor-Hom
identity over $k$, and \eqref{eq:F-left}--\eqref{eq:F-right} give
\begin{align*}
 R\Hom(T_{y,r},T_{d,x})
 &\simeq R\Hom_{S_u}(\ell_{L_u}F_u^{-1}P_r,
                         \mathbb S_{S_u}^{-1}P_x)\\
 &\hspace{1em}\otimes_k
 R\Hom_{S_v}(\mathbb S_{S_v}^{2}P_y,
                         r_{R_v}F_v\mathbb S_{L_v}P_d)\\
 &\simeq \Hom_{S_u}(P_{q_u(r)},P_x)
                  \otimes_k\Hom_{S_v}(P_y,P_{q_v(d)}).
\end{align*}
It is concentrated in degree zero, of dimension one precisely when
$x\leq q_u(r)$ and $q_v(d)\leq y$. These are exactly the mixed
relations of the target, in the projective convention.

The displayed comparison is natural in all four variables and in
the projective maps. It is therefore an isomorphism of the mixed
bimodule with the target incidence bimodule, not just an equality
of dimensions. Choose its basis by the tensor of the two canonical
incidence maps. Composition with maps in either family agrees,
by this naturality, with incidence multiplication. Products inside
each family were already identified by the equivalences; there
are no products of two mixed arrows. Thus, for the sum $T$ of all
the displayed objects, we have the actual algebra isomorphism
\[
 \End(T)^{\mathrm{op}}\cong kJ(vu).
\]
The preceding calculations prove $\Hom(T,T[m])=0$ for $m\ne0$.

Each family generates its entire supported category, because it
is the image of projective generators under equivalences and
external tensor products. These two supported categories generate
the full category: every representation has its upper-supported
subrepresentation and lower-supported quotient, and the same
statement follows for bounded complexes by truncation. Hence $T$
is tilting. Lemma~\ref{lem:tilting-principle} gives the enhanced equivalence. All constructions use finite-dimensional
algebras over $k$, with no characteristic-dependent averaging or
division. This proves the theorem for every field.
\end{proof}

\begin{corollary}\label{cor:cyclic}
Cyclic rotations of the block string of $w$ preserve
$\Db(\Rep_k J(w))$ over every field. In particular
Problem~\ref{prob:swap} holds when $A$ and $B$ are empty.
\end{corollary}
\begin{proof}
Cut the string into the two words before and after the desired
rotation and apply Theorem~\ref{thm:rotation}.
\end{proof}



\section{Failure of block-permutation invariance}
\label{sec:failure}
\label{sec:counterexample}
The obstruction in this section concerns the complete intervals $J(w)$,
not a relative boundary condition or a chosen tilting construction.
All calculations use integers; the auxiliary prime is unrelated to the
characteristic of the representation field.

\begin{lemma}[Integral Euler obstruction]\label{lem:euler-polynomial}
Let $C_P=(1_{x\leq y})$ be the zeta matrix of a finite poset $P$ in a
linear extension. The polynomial
\[
 F_P(t)=\det(tC_P^t-C_P)\in\mathbb Z[t]
\]
is invariant under $k$-linear derived equivalence of its incidence
algebra, for every field $k$.
\end{lemma}
\begin{proof}
The simple classes give an integral basis of $K_0$, and the Euler
matrix in this basis is $E_P=C_P^{-1}$. Indeed, covariant projectives
satisfy $\operatorname{Hom}(P_x,S_y)=k\delta_{xy}$ and have dimension
vectors given by the rows of $C_P$. Incidence algebras have finite
global dimension, so the Euler form is defined on the bounded derived
category. A $k$-linear equivalence induces an integral unimodular
change of basis preserving this form. Consequently
$\det(tE_P-E_P^t)$ is invariant. Finally,
\[
 C_P^t(tE_P-E_P^t)C_P=tC_P^t-C_P,\qquad \det C_P=1.
\]
This proves the claim independently of any convention for the sign
of the Coxeter transformation.
\end{proof}

\begin{theorem}[Computational counterexample]\label{thm:counterexample}
Set
\[
 u=110100,\quad v=11011000,\quad A=\varnothing,\quad B=1100.
\]
These are admissible blocks and exterior strings in
Problem~\ref{prob:swap}. For every field $k$ there is no $k$-linear
triangulated equivalence
\[
 \Db(\Rep_k J(uvB))\simeq\Db(\Rep_k J(vuB)).
\]
In particular the block-permutation assertion of
\cite[Section 9, Conjecture 9.1]{Chapoton} is false.
\end{theorem}
\begin{proof}[Proof by exact finite computation]
The prefix heights of $u$ and $v$ are positive except at their final
position, so both are indecomposable. The complete Dexter intervals
below $1uvB100$ and $1vuB100$ each have 288 vertices and 882 covers.
Their zeta matrices give the following exact residues:
\[
\begin{array}{c|r|r|r}
 w & |J(w)| & \text{covers} & F_{J(w)}(2)\bmod1009\\ \hline
 110100110110001100 & 288 & 882 & 238\\
 110110001101001100 & 288 & 882 & 676
\end{array}
\]
The different residues imply different integer polynomials, and
Lemma~\ref{lem:euler-polynomial} rules out equivalence over every field.

Appendix~\ref{app:certificate} explains the algorithm and checks;
the self-contained implementation and full interval certificates are
provided as supplementary files accompanying the manuscript.
The calculation uses the direct Dexter move rule and does not depend
on any positive equivalence in this paper.
\end{proof}

\begin{remark}
No minimality claim is made.
The counterexample is outside both supplementary positive ranges:
it has three non-elementary blocks, of which two are non-pyramid.
The weaker successes for rotation and special block exchanges therefore
do not imply the original universal assertion.
\end{remark}

\begin{corollary}[Two cyclic orbits, two derived types]\label{cor:two-cyclic-orbits}
Let $u=110100$, $v=11011000$, and $b=1100$. For every field $k$,
the six intervals obtained by permuting these blocks form exactly
two $k$-linear derived equivalence classes:
\[
 \{J(uvb),J(vbu),J(buv)\},\qquad
 \{J(vub),J(ubv),J(bvu)\}.
\]
\end{corollary}
\begin{proof}
Corollary~\ref{cor:cyclic} identifies the three members of each set.
Theorem~\ref{thm:counterexample} separates the displayed representatives
$J(uvb)$ and $J(vub)$. Since the three blocks are distinct, the two
sets exhaust all six permutations. No further computations are needed.
\end{proof}
For this multiset the derived category detects the two cyclic orders.
This statement concerns these six words only.


\section{Abstract equivalence and boundary attachment}
\label{sec:boundary}
The lower truncations can be equivalent without admitting any
boundary-compatible equivalence. We first construct the abstract
equivalence, then show what boundary compatibility would imply,
and finally use the counterexample to rule it out.

\subsection{Rotating the lower truncation}
Put
\[
 W_z=\{e_i(x):x\leq z,\ 1\leq i\leq b(x)\}\subset L_z.
\]
This removes the bottom layer as well as the top-shape layer.
The pullback $Q_{W_z}$ sends $P_x$ to $P_{e_1(x)}$, so is fully
faithful. Let $s_z:\mathcal D_{S_z}\to\mathcal D_{L_z}$ be
extension by zero to the bottom layer. The bottom layer is a lower ideal: a chain index cannot decrease
from a positive value to zero along a cover. Since
$Q_{L_z}P_x=P_{e_0(x)}$, Serre sends it to the injective supported
on bottom vertices over $y\leq x$. This identifies its maps as
well as its support, and gives
\begin{equation}\label{eq:bottom-serre}
 \mathbb S_{L_z}Q_{L_z}\simeq s_z\mathbb S_{S_z}.
\end{equation}
In particular $s_z$ has a right adjoint $s_z^R$, and
\begin{equation}\label{eq:bottom-adjoint}
 s_z^R\mathbb S_{L_z}^2
 \simeq\mathbb S_{S_z}r_{L_z}\mathbb S_{L_z}
 \simeq\mathbb S_{S_z}^2\ell_{L_z}.
\end{equation}
The base Serre functor in these identities is retained in the
mixed Hom calculation below.

\begin{theorem}[Lower-truncation rotation]\label{thm:lower-rotation}
For nonempty Dyck words $u,v$ there is an enhanced equivalence
\[
 \mathcal T_{u,v}:\mathcal D_{L_{vu}}\longrightarrow
                                  \mathcal D_{L_{uv}}.
\]
It has the explicit boundary value
\begin{equation}\label{eq:lower-boundary-image}
 \mathcal T_{u,v}Q_{L_{vu}}
 \simeq\mathcal E_{u,v}
       (\mathrm{id}_{\mathcal D_{S_u}}\boxtimes
                              \mathbb S_{L_v}^2Q_{L_v}),
\end{equation}
where the base factors on the left are identified with
$S_u\times S_v$, and $\mathcal E_{u,v}$ is the padding functor
constructed below. Formula~\eqref{eq:lower-boundary-image} is not
asserted to be a constant pullback for general $u,v$.
\end{theorem}
\begin{proof}
The lower truncation has three consecutive layers
\[
 S_u\times S_v,\qquad W_u\times S_v,\qquad S_u\times W_v.
\]
Indeed, over $(x,y)$ with $a=b(x)$, these correspond respectively
to $i=0$, $1\leq i\leq a$, and $a<i\leq a+b(y)$.
The induced orders and mixed relations follow by restricting
Proposition~\ref{prop:boundary-algebra}; between two different
layers the relation is precisely comparison of the two base
coordinates. The layer unions in this order are lower ideals.

Define a monotone map
\[
 \rho:L_{uv}\longrightarrow S_u\times L_v,
 \qquad
 e_i(xy)\longmapsto
 \begin{cases}
 (x,e_0(y)),&i\leq b(x),\\
 (x,e_{i-b(x)}(y)),&i>b(x).
 \end{cases}
\]
The cover rule proves monotonicity. Let
$\mathcal E_{u,v}=\rho^*$, which is exact.
It sends $P_x\boxtimes P_{e_0(y)}$ to the bottom projective of
$L_{uv}$ over $(x,y)$, and sends $P_x\boxtimes P_b$, $b\in W_v$,
to the projective in its last layer. These are actual representable
maps, so $\mathcal E_{u,v}$ is derived fully faithful. It retains
both the first and third layers, filling the middle layer by
pullback of the first-layer coefficients.

Let $a_0:\mathcal D_{W_u\times S_v}\to\mathcal D_{L_{uv}}$
be extension by zero to the middle layer. It is derived fully
faithful. More explicitly, for a middle projective
$P_a\boxtimes P_y$ there is an exact sequence
\begin{equation}\label{eq:middle-resolution}
 0\longrightarrow i_B(P_{q_u(a)}\boxtimes Q_{W_v}P_y)
 \longrightarrow P^{L_{uv}}_{(a,y)}
 \longrightarrow a_0(P_a\boxtimes P_y)\longrightarrow0,
\end{equation}
where $i_B$ is extension from the last layer. Both terms on the
left are projective in the full poset. Their Hom into another
middle-supported object computes exactly the middle Hom; this
also proves the full-faithfulness assertion.

There are no derived maps from the image of $\mathcal E_{u,v}$
to the image of $a_0$: the displayed projective generators of the
first image evaluate to zero on a middle-supported object.
The reverse Hom has a useful precise form:
\begin{equation}\label{eq:three-layer-hom}
 \begin{split}
 R\Hom\bigl(a_0(P_a\boxtimes Y),
               \mathcal E_{u,v}(P_x\boxtimes Z)\bigr)
 \simeq \Hom_{S_u}(P_{q_u(a)},P_x)
                         \otimes_k R\Hom_{L_v}(s_vY,Z).
 \end{split}
\end{equation}
To verify this with its maps, first take $Y=P_y$.
Applying Hom to \eqref{eq:middle-resolution} gives the fibre of
\[
 \Hom_{S_u}(P_{q_u(a)},P_x)\otimes Z(e_0(y))
 \longrightarrow
 \Hom_{S_u}(P_{q_u(a)},P_x)\otimes
                 R\Hom_{W_v}(Q_{W_v}P_y,Z|_{W_v}).
\]
The map is the actual attachment of $Z$ from its bottom layer.
The exact sequence
$0\to i_{W_v}Q_{W_v}P_y\to Q_{L_v}P_y\to s_vP_y\to0$
identifies this fibre with the right side of
\eqref{eq:three-layer-hom}. This argument is natural in all
projective maps and extends to bounded $Y,Z$ by projective
resolutions.

The target $L_{vu}$ has a lower piece $L_v\times S_u$ and an
upper piece $S_v\times W_u$. For their vertices define
\[
 T_{d,x}=\mathcal E_{u,v}
                 (P_x\boxtimes\mathbb S_{L_v}^2P_d),\qquad
 T_{y,a}=a_0(P_a\boxtimes\mathbb S_{S_v}^2P_y).
\]
Internal Hom complexes are those of the two target projective
families. The Hom complexes in one mixed direction vanish as
above. In the other direction,
\eqref{eq:three-layer-hom} and \eqref{eq:bottom-adjoint} identify
$R\Hom(T_{y,a},T_{d,x})$ with
\[
 \Hom_{S_u}(P_{q_u(a)},P_x)\otimes_k
                              \Hom_{S_v}(P_y,P_{q_v(d)}).
\]
It is concentrated in degree zero, with exactly the target
incidence condition. Naturality of the fibre calculation identifies
both module actions and their compositions. Together with the
internal algebra identifications, this proves
$\End(T)^{\mathrm{op}}\cong kL_{vu}$, with no unverified scalar
multiplication choices.

The two families generate the images of $\mathcal E_{u,v}$ and
$a_0$. The first image contains the bottom and last-layer
projectives, and \eqref{eq:middle-resolution} recovers all middle
projectives. Thus the families generate the full source category
and form a tilting object. The resulting equivalence sends target
projectives to the indicated summands. Applying this to the target
bottom projectives gives \eqref{eq:lower-boundary-image}.
More explicitly, a base projective $P_x\boxtimes P_y$ is sent to
$\mathcal E_{u,v}(P_x\boxtimes\mathbb S_{L_v}^2P_{e_0(y)})$.
For $x\leq x'$ and $y\leq y'$, the map from the primed projective
to the unprimed one is the image under $\mathcal E_{u,v}$ of the
representable map in the first factor tensored with
$\mathbb S_{L_v}^2(P_{e_0(y')}\to P_{e_0(y)})$.
These are exactly the internal maps of the first tilting family;
the mixed bimodule identification above is natural for these maps
on both sides. Thus no separate scalar choices are made on the
base arrows. The identification is one of bimodules on the base
projective category and therefore extends to the claimed enhanced
functor identity by bounded projective resolutions.
\end{proof}


\subsection{Extending a boundary-compatible equivalence}
\begin{proposition}[One-sided reduction]\label{prop:one-sided}
For fixed nonempty Dyck words $u,v$, suppose that
$\mathcal D_{J(uvC)}\simeq\mathcal D_{J(vuC)}$ for every nonempty
block string $C$. Then the adjacent exchange holds with all
possibly empty exterior strings $A,B$.
\end{proposition}
\begin{proof}
Put $C=BA$. Theorem~\ref{thm:rotation}, or the identity if $A$ is
empty, gives
\[
 \mathcal D_{J(AuvB)}\simeq\mathcal D_{J(uvBA)},\qquad
 \mathcal D_{J(AvuB)}\simeq\mathcal D_{J(vuBA)}.
\]
Apply the assumed suffix comparison when $BA$ is nonempty.
When $BA$ is empty use Theorem~\ref{thm:rotation} itself. No
compatibility between these rotations is required for existence
of the resulting equivalence.
\end{proof}

The required suffix bimodule has an exact description. More generally,
let $w,w'$ be nonempty Dyck words with a specified poset identification
$S=I(w)\cong I(w')$. For an exchanged pair this is the product-factor
permutation. All functors to and from $S$ below use the specified
identification. For a nonempty suffix $C$, put $V=I(C)$.
By Proposition~\ref{prop:boundary-algebra}, the lower and upper
pieces of $J(wC)$ are
\[
 D_w=L_w\times V,\qquad U=S\times R_C.
\]
The upper piece is literally common to the two words after the
factor identification. The actual mixed derived tensor functor is
\begin{equation}\label{eq:suffix-bimodule}
 \Phi_w=\ell_{L_w}\boxtimes Q_{R_C}:
 \mathcal D_{D_w}\longrightarrow\mathcal D_U.
\end{equation}
Thus replacing $L_w$ by a derived-equivalent category without
controlling $\ell_{L_w}$ does not identify the gluing algebra.
The following criterion transports this specified gluing bimodule.
It uses triangular-algebra tilting, as in
\cite[Theorems 4.5 and 4.9]{LadkaniTriangular}, with the natural
boundary comparison and the Dexter target identification checked here.

\begin{proposition}[Relative extension criterion]\label{prop:relative-extension}
Suppose an enhanced equivalence
$G:\mathcal D_{L_w}\to\mathcal D_{L_{w'}}$ has been constructed,
an enhanced autoequivalence $K$ of $\mathcal D_S$, and a natural
isomorphism of functors
\begin{equation}\label{eq:relative-gate}
 GQ_{L_w}\simeq Q_{L_{w'}}K:\mathcal D_S\longrightarrow
                                      \mathcal D_{L_{w'}}.
\end{equation}
Then for every nonempty $C$ there is a tilting object in
$\mathcal D_{J(wC)}$ with opposite endomorphism algebra $kJ(w'C)$.
The same comparison gives $\mathcal D_{J(AwB)}\simeq
\mathcal D_{J(Aw'B)}$ for arbitrary exterior block strings $A,B$.
\end{proposition}
\begin{proof}
Taking left adjoints in \eqref{eq:relative-gate} gives a natural
isomorphism
\[
 \ell_{L_w}G^{-1}\simeq K^{-1}\ell_{L_{w'}}.
\]
It identifies the mixed functors, with their two module actions:
\begin{equation}\label{eq:suffix-action}
 \Phi_w(G^{-1}\boxtimes\mathrm{id})\simeq
 (K^{-1}\boxtimes\mathrm{id})\Phi_{w'}.
\end{equation}
Let $j_!$ be derived extension from the lower piece $D_w$, and let
$i$ be extension by zero from the upper piece $U$ in $J(wC)$.
For target lower vertices $(d,y)\in L_{w'}\times V$ and target
upper vertices $(x,r)\in S\times R_C$, define
\[
 T_{d,y}=j_!(G^{-1}P_d\boxtimes P_y),\qquad
 T_{x,r}=i(K^{-1}P_x\boxtimes P_r).
\]
Here $j_!$ is essential: replacing it by lower-supported extension
by zero would change the mixed Hom and its degree.
The lower restriction of $j_!E$ is $E$, and its upper restriction
is $\Phi_wE$. Hence $j_!$ and $i$ are fully faithful, and
\[
 R\Hom(j_!E,iW)=0,\qquad
 R\Hom(iW,j_!E)\simeq R\Hom_U(W,\Phi_wE).
\]
Using \eqref{eq:suffix-action}, the remaining mixed calculation is
\begin{align*}
 R\Hom(T_{x,r},T_{d,y})
 &\simeq R\Hom_U(K^{-1}P_x\boxtimes P_r,
                         K^{-1}P_{q_{w'}(d)}\boxtimes Q_{R_C}P_y)\\
 &\simeq \Hom_S(P_x,P_{q_{w'}(d)})\otimes_k
                         \Hom_{R_C}(P_r,Q_{R_C}P_y).
\end{align*}
It is concentrated in degree zero, with value $k$ precisely when
$q_{w'}(d)\leq x$ and $y\leq q_C(r)$. This is the target mixed
incidence relation in the projective convention. Internal Hom
complexes are those of the target projectives by full faithfulness.
Naturality of \eqref{eq:suffix-action} identifies composition with
both internal algebras. There are no products of two mixed arrows,
so these comparisons give $\End(T)^{\mathrm{op}}\cong kJ(w'C)$
including multiplication. All nonzero-degree Hom groups vanish.

The two families generate the images of $j_!$ and $i$. The triangle
$i\Phi_wE\to j_!E\to j_0E$ then recovers every lower-supported
object; lower- and upper-supported objects generate the full
bounded category. Thus $T$ is tilting. For the empty suffix, use
the lower part $L_w$ and the top layer $S$ of $J(w)$: the same
calculation has mixed functor $\ell_{L_w}$ and no suffix tensor
factor. Finally rotate $AwB$ to $wBA$, and $Aw'B$ to $w'BA$,
by Theorem~\ref{thm:rotation}; omit this rotation when $A$ is empty.
The suffix result gives the final assertion.
\end{proof}

\begin{remark}[Boundary data are essential]
The extension criterion is a sufficient theorem with an explicit
mixed bimodule, not an assumption imposed to obtain a universal
exchange. Corollary~\ref{cor:absolute-boundary-obstruction}
shows that its hypothesis fails for every possible pair $(G,K)$
in the counterexample. Equality on $K_0$, or separate isomorphisms
on the values of $GQ$, cannot replace the natural comparison of
both module actions in \eqref{eq:relative-gate}.
\end{remark}


\subsection{The obstruction for every base correction}
\begin{corollary}[Abstract equivalence does not preserve boundary type]
\label{cor:absolute-boundary-obstruction}
Let $u=110100$, $v=11011000$, and identify
$S=I(uv)\cong I(vu)$ by permutation of the two product factors.
For every field $k$ there is an enhanced equivalence
$\mathcal D_{L_{uv}}\simeq\mathcal D_{L_{vu}}$.
Nevertheless there is no pair consisting of an enhanced equivalence
$G:\mathcal D_{L_{uv}}\to\mathcal D_{L_{vu}}$ and an enhanced
autoequivalence $K:\mathcal D_S\to\mathcal D_S$ with a natural isomorphism
\[
 GQ_{L_{uv}}\simeq Q_{L_{vu}}K.
\]
\end{corollary}
\begin{proof}
The inverse of Theorem~\ref{thm:lower-rotation} gives the abstract
equivalence. If the stated pair existed, Proposition~\ref{prop:relative-extension}
with suffix $1100$ would give an equivalence between $J(uv1100)$
and $J(vu1100)$, contradicting Theorem~\ref{thm:counterexample}.
The quantifier on $K$ includes all enhanced base autoequivalences,
not merely products of Serre powers. The enhancement and naturality
requirements are exactly those of the extension criterion.
\end{proof}



\section{Positive supplementary families}
\label{sec:positive}
\subsection{Exchanges involving an elementary block}
The boundary value in Theorem~\ref{thm:lower-rotation} can be corrected when one factor
is the elementary block $t=10$. This is a genuine exterior-string
case, not only a rotation of the whole word.

\begin{theorem}[An elementary block with arbitrary exterior strings]
\label{thm:elementary-context}
Let $t=10$. For every nonempty Dyck word $u$, every pair of possibly
empty block strings $A,B$, and every field $k$, there is an equivalence
\[
 \mathcal D_{J(AutB)}\simeq\mathcal D_{J(AtuB)}.
\]
\end{theorem}
\begin{proof}
The base $S_t$ is a point and $L_t$ is the two-point chain $0<1$.
Its bottom projective $P_0=Q_{L_t}k$ satisfies
\begin{equation}\label{eq:two-chain-serre}
 \mathbb S_{L_t}^2P_0\simeq P_1[1].
\end{equation}
For an explicit check, $\mathbb S P_1=P_0$ and
$\mathbb S P_0$ is the bottom simple; applying Serre to the
projective resolution of that simple gives \eqref{eq:two-chain-serre}.

In $L_{ut}$ the last layer has one point $e_{b(x)+1}(xt)$ over
each $x\in S_u$. It is a monotone top section, and every point
projecting below $x$ lies below this section over $x$.
Let $Z_u:\mathcal D_{S_u}\to\mathcal D_{L_{ut}}$ be extension
from this top section. Its projectives are those of the last
layer. Top injectives therefore give, with their actual maps,
\begin{equation}\label{eq:elementary-top-serre}
 \mathbb S_{L_{ut}}Z_u\simeq Q_{L_{ut}}\mathbb S_{S_u}.
\end{equation}
There is no additional top-shape layer in this assertion: it is
an identity inside the lower truncation $L_{ut}$.

Equations \eqref{eq:lower-boundary-image} and
\eqref{eq:two-chain-serre} identify
\[
 \mathcal T_{u,t}Q_{L_{tu}}\simeq Z_u[1].
\]
Consequently the actual equivalence
\[
 H_u=\mathbb S_{L_{ut}}\mathcal T_{u,t}[-1]:
       \mathcal D_{L_{tu}}\longrightarrow\mathcal D_{L_{ut}}
\]
satisfies $H_uQ_{L_{tu}}\simeq Q_{L_{ut}}\mathbb S_{S_u}$.
Its inverse $G_u=H_u^{-1}$ satisfies the relative comparison
\[
 G_u Q_{L_{ut}}\simeq Q_{L_{tu}}\mathbb S_{S_u}^{-1}.
\]
These are natural isomorphisms, obtained on base projectives from
representable and injective maps and then extended through the
specified derived functors. Thus
Proposition~\ref{prop:relative-extension} applies with
$K=\mathbb S_{S_u}^{-1}$. It constructs the suffix tilting object
and identifies its endomorphism algebra for every nonempty suffix.
The empty-suffix case is Theorem~\ref{thm:rotation}.
Proposition~\ref{prop:one-sided} now supplies all $A,B$.
No step uses division by the characteristic of $k$.
\end{proof}

\begin{corollary}\label{cor:two-non-elementary}
If a Dyck word has at most two blocks different from $10$, all
permutations of its blocks give derived-equivalent $J$ intervals
over every field.
\end{corollary}
\begin{proof}
Use Theorem~\ref{thm:elementary-context} to move every elementary
block to the right. With zero or one other block this gives one
normal form. With two other blocks $u,v$ the two possible forms
are $uvt^m$ and $vut^m$. Rotate the first to $vt^mu$ using
Theorem~\ref{thm:rotation}, and move the elementary blocks across
$u$ using Theorem~\ref{thm:elementary-context}. This gives the
second form, including the case $m=0$.
\end{proof}

\begin{remark}[Scope]
The elementary construction supplies an actual pair $(G,K)$.
Corollary~\ref{cor:absolute-boundary-obstruction} proves that no such
pair can exist uniformly for general blocks. The next result
establishes the additional pyramid family without asserting any
extension to all indecomposable blocks.
\end{remark}


\subsection{Pyramid transfers and exchanges}
Write $p_m=1^m0^m$ for $m\geq1$; in particular $p_1=10$.

\begin{theorem}[Pyramid transfer]\label{thm:pyramid-transfer}
For $a\geq1$, $b\geq2$, every field $k$, and arbitrary exterior
block strings $A_0,B_0$, there is an equivalence
\[
 \mathcal D_{J(A_0p_a p_b B_0)}\simeq
 \mathcal D_{J(A_0p_{a+1}p_{b-1}B_0)}.
\]
\end{theorem}
\begin{proof}
Appendix~\ref{app:pyramid-application} identifies the two actual lower
truncations with the posets of Theorem~\ref{thm:relative-mutation}
and proves the required natural boundary comparison. Applying
Proposition~\ref{prop:relative-extension} then gives the assertion
in every exterior context.
\end{proof}

\begin{corollary}\label{cor:pyramid-swap}
Any two pyramid blocks can be exchanged in arbitrary exterior
strings, over every field. Consequently all block permutations
are covered for words with at most one non-pyramid block.
\end{corollary}
\begin{proof}
Successive transfers identify both pairs $(p_a,p_b)$ and
$(p_b,p_a)$ with $(p_{a+b-1},p_1)$ in the same exterior context;
a pair whose second index is one already has that form.
For the last assertion, use two-factor rotation to place the
possible non-pyramid block first. The remaining pyramid blocks
can be permuted by the proved adjacent exchanges, and rotation
places the exceptional block at any desired position.
\end{proof}

\begin{remark}[Scope of the new construction]
This closes an infinite family containing arbitrarily many
non-elementary blocks and arbitrary exterior strings. It does
not settle arbitrary pairs of indecomposable Dyck blocks: their
lower intervals need not admit the Boolean marked-cut model used
in the transfer proof. No such model is assumed for a general
block. The unrestricted Problem~\ref{prob:swap} is false by
Theorem~\ref{thm:counterexample}; that example has two non-pyramid
blocks and lies outside this construction.
\end{remark}



\appendix
\section{The exact computational certificate}
\label{app:certificate}
This appendix specifies the finite calculation used in
Theorem~\ref{thm:counterexample}. It requires no chain-coordinate
formula and no tilting construction.
Generate all Dyck words of semilength 11 by recursively appending $1$
when fewer than 11 have been used, and $0$ when the current height is
positive. Match each opening step to its closing step by a stack.
For every matched substring $s$, if its preceding letter is $0$ and
its following letter is $1$ or absent, move $s$ left past each possible
positive number of consecutive preceding zeros. These are precisely
the Dexter covers in the definition. Store the resulting incoming
edges, and take the reverse reachable set from each of the two
specified upper words. This finite graph has 58786 vertices before
restricting to either interval; the two restricted cover lists are
part of the certificate.

An upward move strictly increases the binary word lexicographically.
Thus sorting the selected vertices gives a linear extension. Starting
with each vertex itself, accumulate its upward reachable set in
reverse lexicographic order; this produces $C$. Evaluate $\det(2C^t-C)$
by integer fraction-free elimination. With previous pivot $d$ (initially
1) and current pivot $a_{kk}$, each trailing entry is replaced by
\[
 a_{ij}\longleftarrow
 \frac{a_{kk}a_{ij}-a_{ik}a_{kj}}{d}\qquad(i,j>k).
\]
Use a row interchange when necessary, record its sign, and take the
last diagonal entry with that sign. Every division is checked exact.
Reducing the resulting integers modulo 1009 gives the values
in Theorem~\ref{thm:counterexample}.

\subsection*{Supplementary files and reproduction}
The following two files are supplied with the manuscript in the
\texttt{supplementary/} directory:
\begin{itemize}
\item \texttt{verify\_counterexample\_independent.py}, the self-contained
implementation of the direct move algorithm and integer elimination;
\item \texttt{independent\_counterexample\_certificate.json}, containing
both full vertex and cover lists, both integer determinants, and the
published control results.
\end{itemize}
The accompanying \texttt{README.md} records the run instructions.
The implementation was tested with Python 3.12.10 and uses only the
standard-library modules \texttt{json} and \texttt{pathlib}.
From the manuscript directory, run
\begin{verbatim}
python supplementary/verify_counterexample_independent.py
\end{verbatim}
Equivalently, from the \texttt{supplementary/} directory, run
\begin{verbatim}
python verify_counterexample_independent.py
\end{verbatim}
Run Python with assertions enabled, without the \texttt{-O} option.
The script writes the certificate JSON beside itself and overwrites
that file on each run, independently of the working directory.
The expected output, in the order
$(\text{vertices},\text{covers},F(2)\bmod1009)$, is
\[
\begin{array}{c|c}
 w & \text{expected tuple}\\\hline
 110100110110001100 & (288,882,238)\\
 110110001101001100 & (288,882,676)
\end{array}
\]
The two control intervals each have 27 vertices and
$F(2)=95869805$. The implementation imports none of the
chain-coordinate implementations used elsewhere in this paper.

\subsection*{Checks on the computation}
The computed vertex and cover sets agree exactly with those from
reverse moves and with Proposition~\ref{prop:covers}. An independent
modular elimination gives the same residues. As a published control,
for $J(11011000)$ and $J(11101000)$ this direct implementation gives
27 vertices and $F(2)=95869805$ in both cases, agreeing with the Coxeter
polynomial $\Phi_2\Phi_4\Phi_{18}\Phi_{54}$ stated in
\cite[Section 9]{Chapoton}, after the change from Coxeter to Serre sign.
The finite computation is an explicit part of this proof, rather than
a proposed symbolic identification of the characteristic polynomials.

\section{The relative mutation underlying pyramid transfer}
\label{app:mutation}
This construction changes a specified family of cross relations.
Its input is a specified finite poset, not a presumed relative
equivalence. It will apply to a transfer between two pyramid blocks.

Let $R$ be a finite poset with a least element. Let $A,B$ be finite
posets with monotone maps $q_A,q_B$ to $R$. Assume that $A$ has a bottom copy
$\alpha(R)$ and that $B$ has a bottom copy of $R$, both lower ideals,
with
\[
 \alpha(r)\leq a\ \Longleftrightarrow\ r\leq q_A(a),\qquad
 r\leq b\ \Longleftrightarrow\ r\leq q_B(b),
 \qquad q_A\alpha=q_B|_R=\mathrm{id}_R.
\]
Put $V=B\setminus R$. Form
\[
 D=A^0\sqcup\widehat R\sqcup B^0,\qquad
 U=A^1\sqcup B^1.
\]
Within $D$ and $U$, retain the given internal orders and compare
vertices from an earlier displayed piece to a later displayed
piece by comparison of their images in $R$, with no relations in
the reverse direction. Define $P$ with lower ideal $D$ and upper filter $U$
by the following mixed relations, and no others:
\begin{align*}
 a^0\leq (a')^1&\ \Longleftrightarrow\ a\leq a',&
 a^0\leq b^1&\ \Longleftrightarrow\ q_A(a)\leq q_B(b),\\
 \widehat r\leq b^1&\ \Longleftrightarrow\ r\leq q_B(b),&
 b^0\leq (b')^1&\ \Longleftrightarrow\ b\leq b'\text{ and }b'\in V.
\end{align*}
Let $P'$ have the same pieces and relations except that its last
condition is $b\leq b'$, without the requirement $b'\in V$.
These relations are transitive: internal comparisons use monotonicity
of $q_A,q_B$, and the restricted last mixed relation is stable under
postcomposition because $V$ is an upper filter in $B$. Precomposition
from $A^0$ or $\widehat R$ is covered by the other displayed mixed
relations. No relation points backwards across the ordered pieces,
so antisymmetry also holds. Thus both constructions are posets.
Projection gives monotone maps
$P,P'\to R\times C_2$, where $C_2=\{0<1\}$ records $D$ or $U$.
Write $Q_P,Q_{P'}$ for their pullbacks.

\begin{theorem}[Relative two-term mutation]\label{thm:relative-mutation}
There is an enhanced equivalence $G:\mathcal D_P\to\mathcal D_{P'}$
with a natural identification
\begin{equation}\label{eq:mutation-boundary}
 GQ_P\simeq Q_{P'}
       (\mathrm{id}_{\mathcal D_R}\boxtimes\mathbb S_{C_2})[-1].
\end{equation}
This construction works over every field.
\end{theorem}
\begin{proof}
All projectives in the construction below belong to $P$.
For $d\in A^0\sqcup\widehat R\sqcup V^0$ define
\[
 f(a^0)=a^1,\qquad f(\widehat r)=r^1,\qquad f(v^0)=v^1,
 \qquad C_d=\operatorname{Cone}(P_{f(d)}\longrightarrow P_d).
\]
The arrow is the incidence inclusion. Its restriction to $U$ is
an isomorphism onto $P_d|_U$, so $C_d$ is the extension by zero of
the projective at $d$ in $D$. For $r\in R$ put
\[
 F_r=\operatorname{Cone}
   \bigl(P_{r^0}\oplus P_{r^1}
                      \xrightarrow{(1,1)} P_{\widehat r}\bigr).
\]
Thus each displayed cone is a projective complex in degrees $-1,0$.
There is a natural triangle
\begin{equation}\label{eq:mutation-F-triangle}
 P_{r^0}\longrightarrow C_{\widehat r}\longrightarrow F_r
                              \longrightarrow P_{r^0}[1].
\end{equation}
Index a family of objects by the vertices of $P'$ as follows:
\[
\begin{array}{c|ccccccc}
 \text{vertex}&a^0&\widehat r&r^0&v^0&a^1&r^1&v^1\\\hline
 T_{\text{vertex}}&P_{a^1}[1]&P_{r^1}[1]&F_r&P_{v^1}[1]
                    &C_{a^0}&C_{\widehat r}&C_{v^0}
\end{array}
\qquad(v\in V).
\]
We check its entire Hom algebra, including the new $F_r$ terms.
Let $j_0$ and $i$ be the lower- and upper-supported embeddings.
For the indicated $d,e$ and $u,t\in U$, the projective resolutions
above give
\begin{align*}
 R\Hom(C_d,C_e)&=\Hom_D(P_d,P_e),\\
 R\Hom(P_u[1],P_t[1])&=\Hom_U(P_u,P_t),\\
 R\Hom(P_u[1],C_d)&=0,\\
 R\Hom(C_d,P_u[1])&=\Hom_U(P_{f(d)},P_u).
\end{align*}
All complexes on the right are in degree zero. These comparisons
include the usual incidence maps.

For the remaining terms, direct application of the same two-term
resolution gives
\[
 R\Hom(C_d,P_{r^0})=0,\qquad
 R\Hom(C_{\widehat r},C_d)\xrightarrow{\sim}
                                  R\Hom(P_{r^0},C_d).
\]
For $d\in V^0$, the first vanishing is the cancellation of two
identical incidence maps; for $d\in A^0\sqcup\widehat R$ both
terms are zero. In the second identity both terms evaluate the
same lower projective on the bottom of $B$. Applying
\eqref{eq:mutation-F-triangle} now yields
\begin{align*}
 R\Hom(F_r,C_d)&=0,&
 R\Hom(C_d,F_r)&\simeq\Hom_D(P_d,P_{\widehat r}),\\
 R\Hom(F_r,F_s)&\simeq\Hom_R(P_r,P_s),&
 R\Hom(F_r,P_u[1])&\simeq\Hom_U(P_{r^1},P_u).
\end{align*}
In the other direction,
\[
 R\Hom(P_u[1],F_r)\simeq
 \begin{cases}
  k[0],&u=v^1,\ v\in V,\ r\leq q_B(v),\\
  0,&\text{otherwise}.
 \end{cases}
\]
Indeed $\Hom(P_u,P_{r^0})$ is exactly the missing-bottom mixed
bimodule, and the two shifts in the triangle cancel.
These identities prove, for every pair of target vertices,
\[
 R\Hom(T_x,T_y)=
 \begin{cases} k[0],&y\leq x\text{ in }P',\\0,&\text{otherwise}.
 \end{cases}
\]

Here is also a multiplication check, to exclude a twisted or
zero-relation algebra. Let $a_*=\alpha(\min R)$; the least target
vertex has image $P_{a_*^1}[1]$. Choose maps from every $T_x$ to
this object as follows: use incidence maps on the shifted upper
projectives; on $C_d$ use its degree $-1$ component
$P_{f(d)}\to P_{a_*^1}$; on $F_r$ use the $P_{r^1}$ component.
Whenever $y\leq x$, composition with the chosen map from $T_y$
induces an isomorphism
\[
 \Hom(T_x,T_y)\longrightarrow\Hom(T_x,P_{a_*^1}[1]).
\]
For the families without $F$ this follows from the displayed
projective resolutions. For maps into $F_r$ from a shifted
$P_{v^1}$, the representative has components $(1,-1)$; its
$P_{r^1}$ component is a unit. For maps into $F_r$ from $C_d$,
that same component generates the quotient: homotopies affect
only the $P_{r^0}$ component. The maps between two $F$ terms are
the natural incidence maps of $R$. Maps out of $F$ to an upper
projective are induced by $C_{\widehat r}$, with this component
again a unit. This exhausts the nonzero cases. All assertions use
only cancellation of identity maps and coefficients $1,-1$.
Normalize the generator $a_{yx}$ by requiring its composite to
the least object to equal the chosen map from $T_x$. Associativity
then gives $a_{zy}a_{yx}=a_{zx}$ for $z\leq y\leq x$. Consequently
$\End(T)^{\mathrm{op}}\cong kP'$ including multiplication.

The shifted upper projectives are summands. Their triangles with
$C_d$ recover all $P_d$ for $d\in A^0\sqcup\widehat R\sqcup V^0$.
Triangle~\eqref{eq:mutation-F-triangle} recovers the remaining
$P_{r^0}$. Hence $T$ generates and is tilting. Let
$\mathcal T:\mathcal D_{P'}\to\mathcal D_P$ be the resulting
equivalence, sending the target projectives to this family.

Finally $Q_{P'}(P_r\boxtimes P_0)=P_{\alpha(r)^0}$ and
$Q_{P'}(P_r\boxtimes P_1)=P_{\alpha(r)^1}$. Their images are
respectively $Q_P(P_r\boxtimes P_1)[1]$ and
$Q_P(P_r\boxtimes\operatorname{Cone}(P_1\to P_0))$.
To check that the algebra normalization preserves all these maps,
write
\[
 X_r=P_{\alpha(r)^1}[1],\qquad Y_r=C_{\alpha(r)^0}.
\]
Use the incidence maps on $X_s\to X_r$ and the componentwise
incidence chain maps on $Y_s\to Y_r$ for $r\leq s$.
Let $\delta_r:Y_r\to X_r$ be the connecting map, represented by
the identity on the degree $-1$ term of the cone. The mixed map
$Y_s\to X_r$ is $X_s\to X_r$ composed with $\delta_s$.
It equals $\delta_r$ composed with $Y_s\to Y_r$ as a chain map.

Let $h_Z:Z\to P_{a_*^1}[1]$ be the maps used for normalization
above. On $X_r$ it is the incidence map; on $Y_r$ it is the same
incidence map on its degree $-1$ term. Hence for each of the
three displayed types of base arrow $b:Z\to Z'$ one has
$h_{Z'}b=h_Z$, with coefficient $+1$. These arrows already are
the normalized incidence generators. No further rescaling of a
base arrow or of an individual summand is needed. Moreover for
any normalized arrow $a:W\to Z$, including an arrow involving
$F_r$, one has $h_{Z'}ba=h_W$; uniqueness of the normalized
generator shows that the mixed composite is also preserved.
This checks the base arrows and their compositions with the full
tilting algebra simultaneously. On $C_2$ the objects $X_r,Y_r$
and the connecting map are exactly the images of the two
projectives and their arrow under $\mathbb S_{C_2}^{-1}[1]$.
Thus the comparison is natural on the base projective
category, and extends as an enhanced functor identity:
\[
 \mathcal T Q_{P'}\simeq
 Q_P(\mathrm{id}\boxtimes\mathbb S_{C_2}^{-1})[1].
\]
Taking $G=\mathcal T^{-1}$ proves \eqref{eq:mutation-boundary}.
\end{proof}

\subsection{Verification for pyramid blocks}
\label{app:pyramid-application}
We now verify Theorem~\ref{thm:pyramid-transfer}, including the
identification of the actual gluing relations.
\begin{proof}[Verification of the pyramid transfer]
The interval $I(p_m)$ is the Boolean poset of compositions of $m$:
a composition is represented by its concatenation of pyramid
blocks, and a cover removes one cut between consecutive parts.
For completeness, a legal downward move on a concatenation of
pyramids replaces one pyramid by two consecutive pyramids. An
inner split would leave the moved subpath followed by a zero,
contrary to the movability condition. Iterating gives only such concatenations;
conversely each removal of a cut is the legal full slide merging
two consecutive pyramids. This describes all lower elements and
covers. In particular, the index rule of
Proposition~\ref{prop:covers} applies using just the numbers of
current blocks, independently of their lengths.

Put $x\in I(p_a)$, $z\in I(p_{b-1})$ and
$R=I(p_a)\times I(p_{b-1})$. On the source base $I(p_a p_b)$,
mark the first internal cut in $p_b$; on the target base
$I(p_{a+1}p_{b-1})$, mark the last internal cut in $p_{a+1}$.
Writing $0$ for an unremoved cut and $1$ for a removed cut,
both bases are $R\times C_2$. At coordinate $0$ their lower word
is exactly $x\,10\,z$. At coordinate $1$, it is respectively
$x\,\lambda(z)$ or $\rho(x)\,z$, where $\lambda$ increases the
first pyramid of $z$ by one and $\rho$ increases the last pyramid
of $x$ by one. Neither operation changes the number of blocks.
The covers within coordinate $1$ therefore identify its lower
truncation with $L_{p_a p_{b-1}}$ on both sides.

Set
\[
 A=L_{p_a}\times I(p_{b-1}),\qquad
 B=I(p_a)\times W_{p_{b-1}}.
\]
Here $A$ has its bottom section and $B$ its index-one bottom
section, as required in Theorem~\ref{thm:relative-mutation}.
The coordinate-zero part is $A^0\sqcup\widehat R\sqcup B^0$:
over $(x,z)$ the three parts have indices
$i\leq b(x)$, $i=b(x)+1$, and $i>b(x)+1$. The coordinate-one
part is $A^1\sqcup B^1$, with its cut after $b(x)$.

We check the mixed relations, including the exceptional index.
On the source, removing the marked cut merges block number
$\ell=b(x)+2$. All of $A^0$ crosses with its index unchanged;
$\widehat R$ crosses to the bottom of $B^1$. On $B^0$, local
index $j>1$ crosses to the same local index in $B^1$, whereas
$j=1$ is the forbidden index $i=\ell$. It must first increase
to $j=2$; if this point does not exist, no crossing from that
fibre exists. Subsequent covers cannot decrease a positive
local index greater than one to one. Thus the mixed relation on
$B^0\times B^1$ is exactly $b_0\leq b_1$ in $B$, with
$b_1\in B\setminus R$. Conversely every such relation is
obtained by crossing directly for $j>1$, or first increasing to
two for $j=1$, and then following the covers of $B$. All other
mixed relations are exactly those in $P$ above, by the marked
boundary formula.

On the target, the removed cut merges block number $b(x)+1$.
Its forbidden point is now $\widehat R$, which can always first
increase to the bottom of $B^0$. Consequently it still reaches
the bottom of $B^1$. Every point of $B^0$ crosses directly with
its local index unchanged. The mixed relation on $B^0\times B^1$
is therefore the full comparison in $B$. The other relations
are unchanged. This identifies the actual lower truncations
with $P$ and $P'$ of Theorem~\ref{thm:relative-mutation}, not just
their fibre counts.

That theorem constructs the relative equivalence, with base
correction $\mathrm{id}_{\mathcal D_R}\boxtimes\mathbb S_{C_2}[-1]$.
The generalized relative extension criterion
(Proposition~\ref{prop:relative-extension}) transports its actual
mixed bimodule and multiplication to arbitrary exterior strings.
\end{proof}

\section{The canonical shape cut and its limitations}
\label{app:shape}
\subsection{The singleton-cut tilting algebra}
Put $P=J(w)$, let $D$ be its lower-shape part and $U$ its top-shape
part, and identify $U=I(w)$. Then $D$ is a lower ideal and $U$ an
upper filter. The map $f:D\to U$, $f(e_i(x))=x$, is monotone, and
\begin{equation}\label{eq:mixed}
 d\leq u\text{ in }P\quad\Longleftrightarrow\quad f(d)\leq u
 \text{ in }U.
\end{equation}
This follows also from \cite[Proposition 5.6]{Chapoton}: the top of
the chain is the least top-shape point above $d$. Thus the mixed
bimodule has values
\[
 M(u,d)=k\quad\text{if }f(d)\leq u,\qquad M(u,d)=0\text{ otherwise},
\]
with all nonzero incidence compositions equal to the corresponding
basis map. In the order $D,U$, the algebra is
$\left(\begin{smallmatrix} kD&0\\M&kU\end{smallmatrix}\right)$.

The underlying equivalence is the singleton flip-flop of
\cite[Corollary 1.3]{LadkaniPosets}. In that result take
$X=D$, $Y=U$, and the order-preserving map $f:D\to U$:
its positive and negative mixed relations are exactly the two
relations below. We retain the explicit tilting calculation to
specify its object and multiplication in our covariant convention.
For any finite posets $D,U$ and monotone $f:D\to U$, let $P_f$ retain
the two internal orders and have mixed relations $d\leq u$ exactly
when $f(d)\leq u$. Let $P_f^{\mathrm{fl}}$ instead have mixed relations
$u\leq d$ exactly when $u\leq f(d)$. These rules define posets.

\begin{proposition}[Explicit singleton-cut tilting]\label{prop:tilt}
For every field $k$, write $P_z=k\Hom_{P_f}(z,-)$. For $d\in D$ put
\[
 C_d=\operatorname{coker}(P_{f(d)}\longrightarrow P_d),\qquad
 T=\bigoplus_{u\in U}P_u\ \oplus\ \bigoplus_{d\in D}C_d[-1].
\]
Then $T$ is a tilting object and
\[
 \End_{\Db(\Rep_k P_f)}(T)^{\mathrm{op}}\cong kP_f^{\mathrm{fl}}.
\]
Consequently the two bounded derived representation categories are
equivalent.
\end{proposition}
\begin{proof}
The indicated map is the identity on the upper part of $P_d$ and
zero on its lower part. Equation~\eqref{eq:mixed}, with $P_f$ in
place of $P$, gives an exact sequence
\begin{equation}\label{eq:resolution}
 0\longrightarrow P_{f(d)}\longrightarrow P_d
      \longrightarrow C_d\longrightarrow0.
\end{equation}
The restriction of $C_d$ to $D$ is the representable at $d$, and
its restriction to $U$ is zero. This is a length-one projective
resolution in the full category, not just in a subcategory.

By Yoneda and this resolution, all the derived Hom groups between
the summands of $T$ are concentrated in degree zero. Their degree-zero
values are
\begin{align*}
 \Hom(P_u,P_v)&=k\quad\Longleftrightarrow\quad v\leq u,\\
 \Hom(C_d[-1],C_e[-1])&=k\quad\Longleftrightarrow\quad e\leq d,\\
 \Hom(C_d[-1],P_u)&=\Ext^1(C_d,P_u)
                  =k\quad\Longleftrightarrow\quad u\leq f(d),\\
 \Hom(P_u,C_d[-1])&=0.
\end{align*}
Unlisted groups are zero. In detail, applying $\Hom(-,P_u)$ to
\eqref{eq:resolution} gives zero in degree zero from $P_d$ and
$P_u(f(d))$ in degree one from $P_{f(d)}$.
Applying $\Hom(-,C_e)$ gives $C_e(d)$ in degree zero and zero in
degree one. Finally $P_u$ is projective and $C_d(u)=0$.
These computations also exclude negative degrees after the stated
shifts. The triangles of \eqref{eq:resolution}, together with the
$P_u$, recover all $P_d$ from $T$. Since a finite poset algebra has
finite global dimension, its projectives generate its bounded
derived category. This proves generation.

It remains to verify multiplication. Use the canonical representable
maps as bases inside $U$ and inside $D$. For the mixed group use
the class of $P_{f(d)}\to P_u$ in the Hom complex of
\eqref{eq:resolution}. If $e\leq d$, the canonical map $C_d\to C_e$
lifts to the diagram of resolutions with terms
$P_d\to P_e$ and $P_{f(d)}\to P_{f(e)}$.
The diagram commutes with coefficient $1$, because $f(e)\leq f(d)$
and all maps are representable maps. Composition with this diagram
and composition with $P_u\to P_v$ therefore send a mixed basis class
to the corresponding basis class, with coefficient $1$ whenever
the endpoints are comparable. There are no composable pairs of
mixed arrows. These checks exhaust the products; identity maps are
preserved. Taking the opposite algebra gives exactly the incidence
basis and multiplication of $P_f^{\mathrm{fl}}$.

The tilting argument of Lemma~\ref{lem:tilting-principle} gives
the equivalence. All entries in
the resolutions and product comparisons are $0$ or $1$, so no
restriction on the field is used.
\end{proof}


\subsection{Why its target is not the exchanged interval}
Applying Proposition~\ref{prop:tilt} to the shape cut constructs a
specific algebra, $kP_f^{\mathrm{fl}}$. Identifying it with
$kJ(AvuB)$ is an additional assertion and can fail for this choice
of cut.

\begin{proposition}\label{prop:wrong-target}
For $w=(10)(1100)$ the flipped shape poset $P_f^{\mathrm{fl}}$ has
two maximal elements. In particular it is not isomorphic to
$J((1100)(10))$, and its incidence algebra is not isomorphic to
that target incidence algebra over any field.
\end{proposition}
\begin{proof}
The lower interval has just $x_0=101010<x_1=101100$. Its sole cover
merges at position $\ell=3$. The lower-shape part consists of chains
\[
 a_0<a_1<a_2<a_3,\qquad b_0<b_1<b_2,
\]
where $a_i=e_i(x_0)$ and $b_i=e_i(x_1)$, with additional generating
relations $a_i<b_i$ for $i=0,1,2$. This follows directly from
Proposition~\ref{prop:covers}. It has exactly two maximal points,
$a_3,b_2$. In the flipped poset all upper-shape points lie below
lower-shape points, so these remain exactly its maximal points.
Every $J(z)$ has a unique maximum. Moreover, for covariant
representations, the simple projective modules of a poset algebra
are exactly those at maximal points: the representable at $z$ is
simple precisely when nothing is strictly above $z$. Thus the two
algebras have different numbers of simple projectives, an
isomorphism invariant over every field.
\end{proof}
This is a failure of the full shape-cut candidate, not a
counterexample to Problem~\ref{prob:swap}. The two nine-element
intervals themselves are known to be derived equivalent by the
smaller flip-flop described in \cite[Section 9]{Chapoton}.
Nor does Proposition~\ref{prop:wrong-target} rule out an additional
tilting step from the displayed flipped algebra.



\section{Obstructions for restricted base corrections}
\label{app:restricted}
\subsection{A strict comparison}
The example in this appendix involves $10$ and $1100$, not the
pair in Corollary~\ref{cor:absolute-boundary-obstruction}.
\begin{proposition}[A necessary base twist]\label{prop:strict-relative-obstruction}
Let $w=(10)(1100)$ and $w'=(1100)(10)$, and identify their lower
intervals $S$ by the product permutation. There is no $k$-linear
triangle equivalence $G:\mathcal D_{L_w}\to\mathcal D_{L_{w'}}$
with $GQ_{L_w}\simeq Q_{L_{w'}}$. Thus setting
$K=\mathrm{id}$ in Proposition~\ref{prop:relative-extension}
is in general too restrictive.
\end{proposition}
\begin{proof}
Both $S$ are the two-point chain $x_0<x_1$, with $x_0=101010$
and $x_1=w$ or $w'$. Order the seven vertices of each $L$ as
$e_0(x_0),\ldots,e_3(x_0),e_0(x_1),\ldots,e_2(x_1)$.
Proposition~\ref{prop:covers} gives their zeta matrices
\[
 C=\begin{pmatrix}
 1&1&1&1&1&1&1\\
 0&1&1&1&0&1&1\\
 0&0&1&1&0&0&1\\
 0&0&0&1&0&0&0\\
 0&0&0&0&1&1&1\\
 0&0&0&0&0&1&1\\
 0&0&0&0&0&0&1
 \end{pmatrix},\qquad C'=C+E_{4,7},
\]
where $E_{4,7}$ is a matrix unit. The two pulled-back base
projectives have the same dimension-vector matrix on both sides:
\[
 Q=\begin{pmatrix}1&0\\1&0\\1&0\\1&0\\1&1\\1&1\\1&1\end{pmatrix}.
\]
For a zeta matrix $C$, the Euler form in the simple basis is
$E=C^{-1}$, and the Serre action is
$\Sigma=E^{-1}E^t=C(C^{-1})^t$. The latter identity follows from
$\chi(a,b)=\chi(b,\Sigma a)$. Therefore the matrix of
$\chi(QP_i,\mathbb S_L^4 QP_j)$ is $Q^t E\Sigma^4Q$.
Direct multiplication of the displayed integer matrices yields
\[
 Q^t C^{-1}\bigl(C(C^{-1})^t\bigr)^4Q
   =\begin{pmatrix}0&-1\\0&-1\end{pmatrix},\qquad
 Q^t(C')^{-1}\bigl(C'((C')^{-1})^t\bigr)^4Q
   =\begin{pmatrix}-1&-1\\0&0\end{pmatrix}.
\]
An equivalence commutes with Serre functors and preserves the
Euler form. The proposed strict relative identification would
therefore identify these matrices entry by entry, a contradiction.
All entries are Euler characteristics over $k$, so this argument
works in every characteristic; the integer invariant is not
reduced modulo the characteristic.
\end{proof}

\begin{remark}
This obstruction does not contradict Theorem~\ref{thm:rotation}
or rule out the extension criterion with a nontrivial $K$.
For example, in the projective basis of the two-point base chain,
its Serre action is
$K_0=\left(\begin{smallmatrix}1&1\\-1&0\end{smallmatrix}\right)$.
The two fourth-moment matrices above satisfy
$M=K_0^t M'K_0$. Thus this particular obstruction disappears when
the base Serre functor is allowed. This equality is a necessary
numerical check, not a construction of $G$ or of its natural
compatibility. Theorem~\ref{thm:elementary-context} constructs
$(G,K)$ when one factor is $10$. For arbitrary factors a uniform
full-context equivalence is impossible by Theorem~\ref{thm:counterexample}.
\end{remark}


\subsection{Products of Serre powers}
The elementary-block construction does not extend merely by choosing
other integer powers of the two base Serre functors. The following
obstruction concerns the sufficient relative criterion, not the
original block-permutation problem.

\begin{proposition}\label{prop:factor-serre-obstruction}
Let $u=1100$, $v=111000$, and identify the bases of $L_{uv}$ and
$L_{vu}$ with $S=S_u\times S_v$. There is no enhanced equivalence
$G:\mathcal D_{L_{uv}}\to\mathcal D_{L_{vu}}$ satisfying
\[
 GQ_{L_{uv}}\simeq Q_{L_{vu}}
       (\mathbb S_{S_u}^{p}\boxtimes\mathbb S_{S_v}^{q})[h]
\]
for any integers $p,q,h$ and any field $k$.
\end{proposition}
\begin{proof}
For $X=L_{uv}$ or $L_{vu}$, let $C_X$ be its zeta matrix in a
linear extension and let $Q_X$ be the dimension-vector matrix of
the pulled-back base projectives. Set
\[
 M_n^X=Q_X^t C_X^{-1}
             \bigl(C_X(C_X^{-1})^t\bigr)^n Q_X.
\]
If the proposed relative equivalence exists and its base action
on the projective Grothendieck basis is $K$, preservation of Euler
pairings and commutation with Serre give the necessary identities
\begin{equation}\label{eq:twisted-moments}
 M_n^{L_{uv}}=K^tM_n^{L_{vu}}K\qquad(n\geq0).
\end{equation}
These are integer Euler characteristics, including in positive
characteristic.

Here $S_u$ is the chain on $1010,1100$, and $S_v$ is the diamond
with ordered vertices
\[
 101010,\quad 101100,\quad 110010,\quad 111000.
\]
In those
projective bases their Serre matrices are
\[
 U=\begin{pmatrix}1&1\\-1&0\end{pmatrix},\qquad
 V=\begin{pmatrix}
 1&1&1&1\\-1&0&-1&0\\-1&-1&0&0\\1&0&0&0
 \end{pmatrix}.
\]
They satisfy $U^3=-I_2$ and $V^3=I_4$. Consequently the congruence
in \eqref{eq:twisted-moments} for
$K=(-1)^h U^p\otimes V^q$ depends only on $p,q$ modulo $3$.

Here is an explicit finite certificate for all nine residues.
Order the base pairs $(x,y)$ lexicographically in the displayed
orders. For each pair use the vertices $e_i(xy)$ on the source
and $e_i(yx)$ on the target, with
$0\leq i\leq b(x)+b(y)$. This gives linear extensions with $36$
vertices on both sides. In these coordinates $Q_X$ has entry $1$
exactly when its column base pair is below its row base pair.
The zeta matrices are obtained from the cover rule of
Proposition~\ref{prop:covers}; thus no choice of mixed maps enters
this computation. Matrix multiplication gives the following
violations of \eqref{eq:twisted-moments}. Entry indices are one-based;
$M_n=M_n^{L_{uv}}$ and $M'_n=M_n^{L_{vu}}$.
\[
\begin{array}{c|c|c|c|r|r}
 p&q&n&(i,j)&(M_n)_{ij}&(K^tM'_nK)_{ij}\\\hline
 0&0&4&(1,4)&-1&0\\
 0&1&4&(1,4)&-1&0\\
 0&2&4&(1,4)&-1&0\\
 1&0&4&(1,4)&-1&0\\
 1&1&4&(1,4)&-1&0\\
 1&2&4&(1,4)&-1&0\\
 2&0&4&(1,1)&0&-1\\
 2&1&5&(1,5)&0&-1\\
 2&2&4&(1,4)&-1&0
\end{array}
\]
Every residue violates a necessary identity. Periodicity and the
irrelevance of the overall sign prove the assertion for all
integer powers and shifts.
\end{proof}

\begin{remark}
The proposition excludes this entire factorwise Serre family,
not arbitrary base autoequivalences. More generally,
\eqref{eq:twisted-moments} implies simultaneous similarity of
$A_n=(M_0)^{-1}M_n$: namely $A_n'=KA_nK^{-1}$.
The matrix $M_0$ is invertible because the pullback is fully faithful
and the base Euler form is nonsingular. Therefore traces of all
words in the $A_n$ are necessary invariants independent of $K$.
Equality of such traces does not identify the corresponding
boundary functors and cannot replace their natural comparison.
No counterexample to Problem~\ref{prob:swap} follows here.
\end{remark}

\begin{remark}[The refined correction is realized]
For $(u,v)=(1100,111000)$ the base is $C_2^3$. Label its
coordinates $(a,b,c)$ with the last two belonging to $S_v$.
The marked coordinate in Theorem~\ref{thm:pyramid-transfer} is $b$;
under the original target factor ordering its relabeling is
$\Pi:(a,b,c)\mapsto(c,a,b)$. Thus that theorem constructs an actual
relative equivalence with base correction
\[
 \Pi\circ(\mathrm{id}\boxtimes\mathbb S_{C_2}
                         \boxtimes\mathrm{id})[-1].
\]
Its matrix is $-P(I_2\otimes U\otimes I_2)$, where $P$ denotes
the coordinate permutation. All the moment identities
\eqref{eq:twisted-moments} now follow from the natural relative
comparison. The construction mixes the internal factors of the
original base, so Proposition~\ref{prop:factor-serre-obstruction}
does not apply to it. An overall shift by two changes $[-1]$ to
$[1]$ without changing that matrix.
\end{remark}



\section*{Acknowledgments}

This mathematical paper was prepared with the assistance of artificial intelligence.  I
first acknowledge the textual data produced by humanity as a whole, which
provides the linguistic and cultural record from which mathematical expression
can be learned.  I then thank the accumulated scholarly achievements of
mathematicians throughout history.  I also thank the researchers whose work has
advanced and refined neural networks through successive generations.  Finally,
I thank the OpenAI team for the training and development of the models used in
the preparation of this paper.

\begin{thebibliography}{9}
\bibitem{Chapoton}
F.~Chapoton, \emph{Some properties of a new partial order on Dyck paths},
Algebraic Combinatorics \textbf{3} (2020), 433--463.
\href{https://doi.org/10.5802/alco.98}{doi:10.5802/alco.98}.
\bibitem{LadkaniPosets}
S.~Ladkani, \emph{Universal derived equivalences of posets},
arXiv:0705.0946v2 (2007).
\href{https://arxiv.org/abs/0705.0946v2}{arXiv:0705.0946v2}.
\bibitem{LadkaniSheaves}
S.~Ladkani, \emph{On derived equivalences of categories of sheaves
over finite posets}, Journal of Pure and Applied Algebra
\textbf{212} (2008), 435--451.
\href{https://doi.org/10.1016/j.jpaa.2007.06.005}{doi:10.1016/j.jpaa.2007.06.005};
\href{https://arxiv.org/abs/math/0610685v1}{arXiv:math/0610685v1},
the version used for the cited numbering.
\bibitem{LadkaniTriangular}
S.~Ladkani, \emph{Derived equivalences of triangular matrix rings
arising from extensions of tilting modules}, arXiv:math/0702060v3 (2008).
\href{https://arxiv.org/abs/math/0702060v3}{arXiv:math/0702060v3}.
\end{thebibliography}
\end{document}
